\documentclass[10pt]{article}

\usepackage[utf8]{inputenc}

\usepackage[T1]{fontenc}

\usepackage{amsmath}

\usepackage{amsfonts}

\usepackage{amssymb}

\usepackage[version=4]{mhchem}

\usepackage{stmaryrd}

\usepackage{graphicx}

\usepackage[export]{adjustbox}

\graphicspath{ {./images/} }

\usepackage{bbold}



\begin{document}

графике (см. Замкнутый график, теорема о замкнутом графике)



Лит.: [1] И о с и д а К., Функциональный анализ, чер. с англ., М., 1967; [2] Р о б е р т с о н А.-П., Р о б е р т с о н В.Д.н., Топологические векторные пространства, пер. с англ., М.,

1967 .



ОТКРЫТОЕ ЯДРО - множество всех точек $x$ подмножества $A$ топологич. пространства $X$, для к-рых существует такое открытое в $X$ множество $U_{x}$, что $x \in U_{x} \subset A$. О. я. множества $A$ обозначается обычно Int $A$ и представляет собої максимальное открытое в $X$ множество, содержащееся в $A$. Имеет место равенство Int $A=X-[X-A]$, где [] обозначает замыкание в пространстве $X$. О. я. множества топологич. пространства $X$ есть открытое р е г у л я р н о е, или каноническое множество. Пространства, в к-рых открытые канонич. множества образуют баз! топологии, наз. п ол у р г у л я р ы и и. Всякое регулярное пространство полурегулярно. Иногда $O$. я. наз. в н у т$\begin{array}{lll}\text { р е н о с т ь ю м н о же с т в а. } & \text { в. И. Пономарев. }\end{array}$



ОТКРЫТО-ЗАМКНУТОЕ МНОЖКСТВО - подмножество топологич. Іространства, одновременно открытое и замкнутое в нем. Топологич. пространство $X$ несвязно тогда и только тогда, когда в нем имеется отличное от $X$ и от $\phi$ О.-3. м. Если семейство всех О.-3. м. топологич. пространства является базой его топологии, то это пространство наз. и н Д у к т и в н о н у л ь м е р н м. Всякая булева алгебра изоморфна булевої алгебре всех 0.-3. м. подходящего индуктивно нульмерного бикомпакта. Особый класс нульмерных бикомпактов образуют т. н. ә к с т р е м а л ь н о нес в щиеся тем, что в них замыкание любого открытого множества также открыто (и замкнуто). Всякая полная булева алгебра изоморфна булевой алгебре всех О.-з. м. подходящего экстремально несвязного бикомпакта.



ОТНОСИТЕЛБНАЯ ГОМОЛОГИЧЕСКАЯ Пономарев. РА - гомологическая алгебра, ассоциированная с па-



\includegraphics[max width=\textwidth, center]{2023_07_08_1841a6686b69fbcab64fg-1}

фұунктором $\Delta: \mathfrak{Y} \rightarrow \mathfrak{M}$. Функтор $\Delta$ предполагается аддитивным, точным и полным. Короткая точная последовательность объектов категории $\mathfrak{A}$



\[

0 \longrightarrow A \longrightarrow B \longrightarrow C \longrightarrow 0

\]



наз. Д о п у с т м о й, если точная последовательHOCTb



\[

0 \longrightarrow \Delta A \longrightarrow \Delta B \longrightarrow \Delta C \longrightarrow 0

\]



расщепляется в категории $\mathfrak{M}$. Посредством класса допустимых точных последовательностей определяется



\includegraphics[max width=\textwidth, center]{2023_07_08_1841a6686b69fbcab64fg-1(3)}

ных) объектов как класс таких объектов $P$ (соответственно $Q)$, для к-рых функтор $\operatorname{Hom}_{\Im}(P,-)$ (соответственно $\left.\operatorname{Hom}_{\Im}(-, Q)\right)$ точен на допустимых короткіх точных последовательностях.



Любої проективный объект $P$ категории $\mathfrak{A}$ яв-



\includegraphics[max width=\textwidth, center]{2023_07_08_1841a6686b69fbcab64fg-1(1)}

в категории $\mathfrak{U}$ достаточно много относительно проективных объектов (т. е. что для любого объекта $A$ из $\mathfrak{Y}$ существует допустимый эпиморфизм $P \rightarrow A$ нек-рого ס्-проективного объекта категории $\mathfrak{A})$. Если в категории $\mathfrak{A}$ достаточно много £-проективных или £-инъективных объектов, то обычные конструкции гомологич. алгебры позволяют строить в этой категории прогзводные функторы, наз. о т н о с и т е л bны м и п о о з в о дым и фу ун тор а ми.



П. р и м е р ы. Пусть $\mathfrak{U}-$ категория $R$-модулей над ассоциативным кольцом $R$ с единицей, $\mathfrak{M}$ - категория множеств, $\quad \Delta: \mathfrak{Y} \rightarrow \mathfrak{M}-ф у н к т о р, \quad$ «абывающий» структуру модуля. В этом случае все точные последовательності допустимы, и в результате получается «абсолютная" (т. е. обычная) гомологіч. алгебра. Если $G-$ группа, то каждый $G$-модуль является, в частности, абелевой группой. Если $R$ является алгеброї над коммутативным кольцом $k$, то каждый $R-м 0^{-}$ дуль является $k$-модулем. Если $R$ и $S$ - кольца и $R \supset S$, то каждый $R$-модуль является $S$-модулем. Во всех этих случаях имеется функтор из одной абелевой категории в другую, определяющий относительные производные функторы.



Лum.: [1] М а к л е й н С., Гомология, пер. с англ., М., 1966; [2] E i le n be r g., M o o r J J. C., Foundations of relative homological algebra, Providence, 1965



ОТНОСИТЕТЬНАЯ МЕТРӤ. Е. Говоров, А. В. Михалев. ограничение рики $\rho$ на подмножество $A$ метрич. пространства $X$, т. е. ограничение отображения $\rho$ квадрата $X \times X$ на квадрат $A \times A \subset X \times X$. Это понятие позволяет рассматривать как метрич. пространство любое его подмно$\begin{array}{ll}\text { жество. } & \text { Б. Пасынков. }\end{array}$



ОТНОСИТЕЛЬНАЯ СИСТЕМА КОРНЕЙ с в я Зн о й редук тивн о й а лге б р ай е ской г р у п пы $G$, определенной над полем $k$,- система $\Phi_{k}(S, G)$ ненулевых весов присоединенного представления максимального $k$-расщепимого тора $S$ группы $G$ в алгебре Ли q этой группы. Сами веса наз. к о р н ями $G$ о т н о си те л ь н о $S$. О. с. к. $\Phi_{k}(S, G)$, рассматриваемая как подмножество своей линейной оболочки $L$ в пространстве $X(S) \bigotimes_{\mathbb{Z}} \mathbb{R}$, где $X(S)-$ группа рациональных характеров тора $S$, является корневой системой. Пусть $N(S)$ - нормализатор, а $Z(S)$ централизатор $S$ в $G$. Тогда $Z(S)$ является связной компонентой единицы группы $N(S)$; конечная группа $W_{k}(S, G)=N(S) / Z(S) \quad$ наз. г р у п по й В е й л я



\includegraphics[max width=\textwidth, center]{2023_07_08_1841a6686b69fbcab64fg-1(4)}

г р у п п о й В е й л я (о. г. В.). Присоединенное гредставление $N(S)$ в $\mathfrak{g}$ определяет линейное представление $W_{k}(S, G)$ в $L$. Это представление является точным и его образ есть Вейля гр уппа системы корней $\Phi_{k}(S, G)$, что позволяет отождествить әти две группы. Ввиду сопряженности над $k$ максимальных $k$-расщепимых торов в $G$ О. с. к. $\Phi_{k}(S, G)$ и о. г. В. $W_{k}(S, G)$ не зависят, с точностью до изоморфизма, от выбора тора $S$ и часто обозначаются просто $\Phi_{k}(G)$ и $W_{k}(G)$. В случае, когда $G$ расщепима над $k, 0$. с. к. и о. г. В. совпадают соответственно с обычной (абсолютной) системої корней и группой Вейля группы $G$. Пусть $g_{\alpha}$ - весовое относительно $S$ подпространство в $\mathfrak{g}$, отвечающее корню $\alpha \in \Phi_{k}(S, G)$. Если $G$ расцепима над $k$, то $\operatorname{dim} g_{\alpha}=1$ для любого $\alpha$ и $\Phi_{k}(G)$ - приведенная спстема корней в общем случае это не так: $\Phi_{k}(G)$ может быть неприведенной, a dim $g_{\alpha}$ может быть больше 1. О. с. к. $\Phi_{k}(G)$ неприводима, если $G$ проста над $k$.



О. с. к. играет важную роль в описании структуры и в классификации полупростых алгебраич. групп над $k$. Пусть $G$ - полупроста и $T$ - максимальный тор, определенный над $k$ и содержащий $S$. Пусть $X(S)$ и $X(T)-$ группы рациональных характеров торов $S$ и $T$ с фиксированными согласованными отношениямі порядка, $\Delta$ - соответствующая система простых корней группы $G$ относительно $T$ и $\Delta_{0}$ - подсистема в $\Delta$, состояшая из характеров, тривиальных на $S$. Пусть также $\Delta_{k}$ - система простых корней в О. с. к. $\Phi_{k}(S, G)$ определенная выбранным в $X(S)$ отношением порядка; она состоит из сужений на $S$ характеров системы $\Delta$. Группа Галуа $\Gamma=\mathrm{Ga}\left(k_{s} / k\right)$ естественно действует на $\Delta$, и набор данных $\left\{\Delta, \Delta_{0}\right.$, действие $\Gamma$ на $\left.\Delta\right\}$ наз. $k$ -



\includegraphics[max width=\textwidth, center]{2023_07_08_1841a6686b69fbcab64fg-1(2)}

Роль $k$-индекса объясняется следующей теоремої: всякая полугростая группа над $k$ однозначно с точностью до $k$-изоморфизма определяется своим классом относительно изоморфизма нац $k_{s}$, своим $k$-индексом и своим аицзотропныля ядрол. О. с. к. $\Phi_{k}(G)$ полностью олрецеляется системой $\Delta_{k}$ п набором таких натуральных чисел $n_{\alpha}, \alpha \in \Delta_{k}$ (равных 1 пли 2), что $n_{\alpha} \propto \in \Phi_{k}(G)$





\end{document}